\origpage{51}{GRAPH MINORS. III}
\stmt
(1.4) separately here because it is much easier to prove than the result of [4], and our proof here yields a very much better estimate of the value of $w$ than that in [4].

The proof of (1.4) occupies the whole paper (except for some remarks at the end) and we begin by sketching it. It is enough to prove (1.4) when $H$ is a ``cylinder'' (a grid with one pair of opposite sides identified) because every planar graph is isomorphic to a minor of some sufficiently large cylinder. We introduce the ``sleeve union'' of two graphs $G_1$, $G_2$, and observe that the tree-width of the resulting graph $G$ is the maximum of the tree-widths of $G_1$, $G_2$, and that $G_1$, $G_2$ are isomorphic to minors of $G$. It follows that it is enough to prove that if a planar graph is not expressible as a sleeve union of two smaller planar graphs, and has no minor isomorphic to the $n \times n$ cylinder, then its tree-width is bounded by a function of $n$. We prove that in any plane drawing of such a graph there do not exist $n + \frac{1}{2} n^2$ vertex-disjoint circuits each inside the next, and deduce that it has tree-width $\le 3(n + \frac{1}{2} n^2) - 2$. These steps are performed more or less in reverse order.

Let us clarify some terminology. A \textit{walk} in a graph $G$ is a sequence $v_1, e_1, v_2, e_2, \dots, v_k, e_k, v_{k+1}$ of vertices and edges of $G$ (where $k \ge 0$), such that for $1 \le j \le k$, the ends of $e_j$ are $v_j$ and $v_{j+1}$. The walk is closed if $v_1 = v_{k+1}$. A \textit{path} is a walk such that $v_j \ne v_{j'}$ for $1 \le j < j' \le k+1$, and a \textit{circuit} is a closed walk with $k \ge 1$ such that $v_j \ne v_{j'}$ for $1 \le j < j' \le k$. We often identify paths and circuits with the corresponding subgraphs of $G$ when no confusion should arise. The \textit{initial} and \textit{terminal} vertices of the walk are $v_1, v_{k+1}$. When $u, v \in V(G)$, a path between $u$ and $v$ is a path with initial vertex $u$ and terminal vertex $v$, and when $X, Y \subseteq V(G)$, a path between $X$ and $Y$ is a path with initial vertex in $X$ and terminal vertex in $Y$. A path \textit{avoids} $Z \subseteq V(G)$ if no vertex of it is in $Z$. When $Z \subseteq V(G)$, $Z$ \textit{separates} $X$ and $Y$ if no path of $G$ between $X$ and $Y$ avoids $Z$. (Here $X, Y$ may be vertices or sets of vertices.) If $G_1, G_2$ are subgraphs of $G$ we denote by $G_1 \cup G_2$ the subgraph of $G$ consisting of the vertices in $V(G_1) \cup V(G_2)$ and the edges in $E(G_1) \cup E(G_2)$. We define $G_1 \cap G_2$ similarly. If $X$ is a vertex or an edge of $G$, or a set of vertices or edges, we denote by $G \setminus X$ the graph obtained by deleting $X$ from $G$.

\sectionhead{2. \textsc{Radius}}

If $G$ is a planar graph, it can be drawn without crossings in the plane, and by a ``drawing'' we shall always mean such a drawing. If $M$ is a drawing of $G$ and $C$ is a circuit of $G$, then every region of $M$ and every vertex and edge of $G$ not in $C$ is either inside or outside $C$ in $M$, in the obvious sense. For every region $R$ of a drawing $M$, we define $d(R)$ to be the minimum value of $k$ such that there is a sequence $R_0, R_1, \dots, R_k$ of regions of $M$, where $R_0$ is the
